\section{Results: Facilitating Strategies and Challenges}
\label{sec:interview-results}

This section reports the thematic analysis of the three transcribed expert interviews.
E1 (interview~1) works in organizational development, strategy, and multi-stakeholder
processes, and combines facilitation with consulting.
E2 (interview~2) maintains internal decision routines and facilitates commissioned workshops.
E3 (interview~3) facilitates large change projects and leadership teams, online and in person.
Names of persons and organizations that appear in the transcripts are not repeated here.
The interviews were coded with the codebook in \autoref{annex:expert-interview-codebook}.
Codes were applied to claims by the interviewee.
One claim could receive several codes.
Interviewer turns, consent, and biographical self-description were left uncoded.
The codebook was fixed before this coding pass.
It had been built from the research questions and from a reading of these same interviews,
so the coding is a structured reading of the material, not a test of a scheme that was blind
to it.
One coder did the coding.
No second coder was used, and nothing in this section is offered as a frequency beyond
these three accounts.
Consecutive rows often continue a single thought.
Those rows are read as one line of argument.

The interviews were asked to serve two purposes.
They describe how a remote group stays with a joint decision, and they are a source of
requirements for an AI moderator.
Both purposes are reported here.
Findings that later fall outside the prototype are included, because they bound what an
in-meeting system can be asked to do.
Design consequences are not settled in this section.

% ---------------------------------------------------------------------------------------------------------------------
\subsection{What has to be true around the decision}
\label{sec:interview-decision-conditions}

All three experts treat the live discussion as dependent on work done before anyone speaks,
and on what happens after the meeting closes.

The first condition is a shared information base (D2).
E1 calls lack of preparation the largest failure of remote meetings: people enter without
having read the document, the group has no basis for a decision, and the meeting is postponed.
She treats the meeting as an investment of organizational time, and she treats unprepared
attendance as a waste of that investment.
E2 makes the same observation from the other side of the format.
People will not enter a workshop unprepared, and they regularly enter a meeting unprepared.
The invitation often contains no information, so participants cannot tell what is to be
decided, what their role is, or why they were invited.
E2 estimates that at least half of the work is preparation: the concrete problem, who must be
present, whether the decision owner is in the room, and which decision procedure will be used.
E3 locates the same work in the briefing with the person who owns the decision.
Goals and non-goals have to be stated there.
Leaders, in E3's account, often have the goal in mind and still fail to say it.
For a short slot with senior decision makers, E3 prepares a decision memorandum: the dilemma,
options including the option of doing nothing, and the costs and benefits of each.
For a team that is meant to decide together, the same briefing asks how far the team is
actually allowed to go.

The second condition is the mandate (D1).
A decision, in these accounts, is not one kind of act.
E1 distinguishes documented formal decisions from informal agreements that rest on trust, and
she says that many of the decisions in her workshops are decisions about next steps rather
than a final choice.
E2 separates reversible decisions, which can be tried under a ``safe enough to try'' rule,
from irreversible ones, which justify a longer and more careful process.
Not deciding is, for E2, itself a decision, because delay has a cost.
E3 asks who has the authority, including budget authority, and who is responsible if the
meeting ends without a decision.
If the team does not agree, E3 puts that responsibility on the leader, who then has to decide
alone.

The procedure is part of the mandate, and it is the point E2 develops furthest.
In E2's account, organizations usually know two procedures, a decision by authority and a
decision by majority.
Majority is described as a poor fit for most organizational choices, because it produces
winners and losers and can exclude a minority whose cooperation the result needs.
E2 prefers consensus for fundamental questions, systemic consensus when several options are
on the table, and a consultative individual decision when one or two people hold the
expertise: they must hear the others, and they then decide and carry the responsibility.
``Disagree and commit'' is named as a practice that keeps dissent from turning into a refusal
to implement.
E1 and E3 describe a related failure in which the procedure is only apparent.
E2's example is a meeting that formalizes a decision already taken.
E3 describes leaders who tell a team that the decision is theirs when it has already been
made, and who want the team only to nod it through.
E3's advice is to say so in advance.
E1 describes hierarchies that have decided and do not want to hear where the decision does
not fit.
In all three accounts, a meeting that asks for commitment without a real scope of decision
damages the result.

The third condition sits outside facilitation (D3).
E2 states it most sharply.
Many meetings exist because decision rights are missing: small choices wait for the next
regular meeting, people who could decide are not allowed to, and the meeting becomes the
symptom of that structure.
Better moderation, on this account, can support a meeting and still leave the structural
problem in place.
The practical test E2 uses is whether the person who can decide is in the room.
If not, the group needs an explicit way to handle that absence, or it will leave without a
decision.
E2 also treats meetings that only broadcast information as a misuse of synchronous time.
E1's version of the same limit is the guest list.
Standard meetings contain people who cannot contribute and who cannot decide.
Hearing something is not, for her, a reason to attend, because a short protocol can do that
afterwards.
She also says that an external facilitator can invite dissent for the duration of a workshop
and cannot change a fear-based culture.
E3's version is the boundary of the single appointment.
A consequential decision has a before and an after: options have to be clarified, and someone
has to implement the result by a stated time.
The decision is written into the system the organization already uses, so that people who
were not in the room can act on it.

Ownership closes this set of conditions (D2, F5).
E1 returns to it at length.
A result that the leadership team does not take as its own is, in her words, waste.
Ownership comes from having worked on the result and from hearing one's own voice in it.
It does not come from receiving a finished proposal from above, or from leaving the content
with a consultant.
E2 states the operational form: who does what by when, and where the group checks that it
was done.
A decision that nobody follows is, for E2, the worst outcome.
E3 separates that ownership from the facilitator's role, which is taken up below.

% ---------------------------------------------------------------------------------------------------------------------
\subsection{What puts the decision at risk while the group is together}
\label{sec:interview-meeting-dynamics}

Inside the meeting, the three accounts converge on three risks: loss of the decision itself,
suppression of relevant dissent, and signals that are hard to read at a distance.

Loss of the decision (M1) takes several forms.
E3 describes side tracks that open because the decision is not yet ripe, even when it was
prepared, and unstructured talk in a group of ten or twelve, during which people start
writing messages and stop listening.
E2 describes drift away from the problem, and a further case in which the divergence itself
shows that the group has been working on the wrong question.
The response is then a judgment, not an automatic return to the agenda: go back, or continue
with the present question for now.
Time pressure is tied to the middle of the process.
E1 describes the open phase of the diamond, when new information has just come back from
small groups, as the critical moment, because that is when the group has to find its way
toward a decision again.
E2 names the groan zone, in the transcript as ``Grownzone''.
The facilitator has to be able to stay in it.
A complex topic in a forty-five-minute meeting does not leave enough time there, and the
pressure to decide then closes the search too early.
E1 makes the same observation without the name: once the content starts, the hour is already
gone, and the alternative is fewer people and more time.
E3 adds a timing rule at the other end.
Reopening a large controversy in the last third of a workshop is dangerous.
The move then is to stop discussing the substance and to discuss only how a decision will
be reached.

Dissent and hierarchy (M2) are described as conditions of decision quality, not as a problem
of politeness.
E1 says that few organizations treat contradiction as useful.
Many are run through fear, and the hierarchy wants execution rather than news of where a
plan does not fit.
Dissent that is not invited stays at the bottom and later returns as conflict.
Trust is, for her, a condition of being able to say that something is wrong, and it is harder
when the people have never met.
She also describes the opposite of useful dissent: chronic objection, which has to be weighed
rather than simply amplified, and the habit of taking a substantive criticism as a personal
attack.
E2's critical moments are participation failures of three kinds: people do not take part,
they are given no way to take part, or they do not want to.
The procedure can produce the same failure.
A majority vote can leave two people as losers of a decision the whole team then has to live
with.
People who will be affected by a decision and who are not in the room are a further gap.
E2's example is a shift pattern decided without the person who has to work it.
E3 defines the relevant voice by what would happen afterwards.
The people who should speak are those who, if the decision goes against them, could
credibly do otherwise later, including in silence.
A vote of confidence is, for E3, the minimum even when a speech is not required.
Online, that confidence is hard to see.
A thinner ``yes'', less enthusiasm as the decision approaches, or growing quiet is a sign
that something is still missing.
E3 refuses a single reading of that sign.
The same word means different things for different people, and in some settings a lower
rank will not agree until the senior person has agreed, while in others the senior person
stays noncommittal until everyone else has spoken.

The remote setting changes which of these signs can be seen (M3).
E1 prefers to do complex content in person.
In a wall of video tiles she loses the sense of mood, resistance, and agreement that she
has when she is in the room.
Camera-off participation and parallel work are, for her, forms of absence.
A black tile is something she does not accept, because the group then cannot tell whether
a person is present.
Language adds a further gap.
A meeting in a shared foreign language, with more than one cultural habit of what may be
said, makes a common picture harder to reach.
E3 describes a loss that is easy to miss: in a face-to-face break, people repair small
misunderstandings informally, and an online break disconnects them, so those repairs do not
happen.
Shared visual tools are, for E3, a condition of a complex decision online, and so is the
time required for people who do not yet know the tool.
Technical failure across organizations, including the ordinary failure of a microphone, is
named as a further remote obstacle.
E3 also treats tone as data.
An enthusiastic yes that already contains the next action is a different commitment from a
bare yes, and that difference is carried by delivery.
E2 is more sparing on the remote medium.
The points that do appear are parallel work while ``listening'', and the dependence of some
methods on a shared board.
The remote-specific material is therefore strongest in E1 and E3.
The material is a list of cues and repairs that become unreliable when the group is not in
one room.

% ---------------------------------------------------------------------------------------------------------------------
\subsection{What the facilitators do in the convergent phase}
\label{sec:interview-facilitation-practice}

The practices are reported here as practices.
They are the material from which requirements can be taken.
They are not yet requirements.

The shared shape of a convergent meeting (F1) is to open before closing, and to keep
clarification separate from debate.
E2 states the sequence as information, then a check that people mean the same thing, then
perspectives, and only then evaluation.
Jumping from a document into debate is named as a source of decisions that people would
have made differently if they had understood the point.
E1 uses a check-in so that every person has spoken once before the extroverted voices take
the floor, and small groups so that people can think by speaking.
She describes the overall shape as a diamond: build the complexity, then bring it back to a
decision or to next steps.
E3 follows a prepared plan and leaves it when the discussion becomes controversial in a way
the plan did not expect.
Early or mid-way, the two sides may go into separate groups to sharpen what they do not
understand in the other position.
Near the end, the structure changes level.
The group stops reworking the controversy and agrees the path to a decision: which queries,
who fetches which information by when, and when they meet again.

Making the state visible (F2) is the convergent move all three describe in concrete terms.
E1 clusters results on a shared board and then summarizes: given what has been heard, can
the group agree these next steps.
The summary is offered as a reading of the group's material.
E2 has people write pros and cons, or the consequences of a candidate decision, so that the
option space is outside any one person's memory.
With more than a few people, writing or a small group is also a way to keep the contributions
in view.
E3's criterion is that everyone should see, at one glance, where everyone stands.
A serial round in which each person says an opinion is, in a larger group, a practice E3
rejects, because people stop listening.
Scales, a fist-to-five that includes an explicit veto, a matrix of effort against value, and
a poll are ways to put that picture in front of the group.
E3 also paraphrases near the close, sometimes slightly past what was said, in order to test
whether something is still unsaid.
The summary and the display are checks.
They are not the decision.

Bringing a voice or a fact in (F3) is selective.
E1 invites people who have not yet spoken, and she prepares questions that are meant to
produce a new angle on the topic.
She also asks leadership, or a product owner, to say what they will do with what they have
heard.
The invitation of a quiet person is, for her, easier in the room than online, where it can
feel intrusive, and it can stigmatize, especially across culture and age.
Naming introversion as the reason is something she would do cautiously.
E2's corresponding rules are to hear lower rank before senior rank, so that the senior view
does not set the room, and to let everyone speak before anyone evaluates.
A role, such as a skeptic's, is a way to make a critical view sayable.
If a needed role is absent, E2 will ask the group what that role would be likely to say, and
will mark the answer as a hypothesis.
E3 will not invite every silent person.
An intern who says nothing is left.
A co-founder who has been silent on a topic that matters to him is asked, by name, after a
short restatement of what has been said.
The people who are asked are those whose later resistance would matter.
Alongside speech, E1 and E2 both use writing and chat, because some people think before they
can speak and others think by speaking.
E3's minimum, when a speech is not needed, is still a visible degree of confidence.

The judgment of when to move (F4) is described as restraint as much as intervention.
E2 has to endure the groan zone and not force a close.
E3 waits, counts about seven seconds, and only then asks in another way or repeats the
paraphrase.
A non-answer is sometimes the information.
E3 stops pushing when further questioning would be delicate.
E1 watches whether the group is only functionally present or also engaged, and she treats
absence of energy as a reason the process will stall.
She also weighs chronic objection rather than inviting it without limit.
E2 asks, before an irreversible decision, whether the moment is even a fit one.
The example is late on a Friday, with no decision capacity left.
The same expert uses the state of the people, including hunger, as a reason to distrust the
moment.
None of the three describes a cue that would justify the same intervention every time.

The boundary of the role (F5) is drawn most sharply by E3.
Leaders sometimes want the facilitator to produce the decision and to make people accept it.
E3 refuses that delegation.
The facilitator's work is to structure the conversation so that there are real options and
so that people can stand behind a result.
Saying yes or no remains the leader's work, because the facilitator is not in that role.
E2 draws the boundary at the other end of the system: the decision should not be taken out
of human hands, particularly where the solution is not yet known.
E1 requires the group to take ownership, and she also describes a practice that crosses the
line E3 draws.
Because her facilitation sits inside consulting, she prepares substantive proposals for how
to continue, and she describes the art of the work as helping to make the decision as well
as advising on it.
The three accounts agree that the result has to become the group's.
They differ on how far the facilitator may put content into that result.
E3 forbids it.
E1 treats a consulting proposal she has prepared as material on which the group then
decides.

% ---------------------------------------------------------------------------------------------------------------------
\subsection{Where an AI moderator would help, and where it would not}
\label{sec:interview-ai-cofacilitator}

The experts were not asked to design a technical architecture.
They were asked which behavior would be useful, how a system should enter the meeting, and
what it should not do.
\autoref{tab:interview-ai-functions} and \autoref{tab:interview-ai-limits} record which of
them stated which point.
An empty cell means that this expert did not state it.
It does not mean that they rejected it.

\begin{table}[htbp]
  \centering
  \caption{Functions the experts described as useful for an AI in or around a convergent
    meeting.}
  \label{tab:interview-ai-functions}
  \footnotesize
  \begin{tabularx}{\textwidth}{@{}>{\RaggedRight\arraybackslash}Xccc@{}}
    \toprule
    Function & E1 & E2 & E3 \\
    \midrule
    Fast summary, clustering, or a visual restatement of options & yes &     & yes \\
    Timekeeping, agenda position, and a prompt to come to the point & yes &     &     \\
    Walking the group through a decision: question, options, doing nothing,
      pros and cons, or a named procedure &     & yes & yes \\
    Pointing at a decision the group has not noticed, or at agreement that is
      thinner than it appears &     &     & yes \\
    A cue that a named bias may be in play &     & yes &     \\
    A persona that speaks for an absent stakeholder or a missing expertise & yes & yes &     \\
    Documentation into a knowledge base, a task system, or an approved summary & yes &     & yes \\
    Help with preparation, or with stating the goal at the start & yes & yes & yes \\
    \bottomrule
  \end{tabularx}
\end{table}

E1 considers transcription and summary already available, and she considers faster
sensemaking the use for which an AI is well suited.
Clustering a board that would otherwise remain a mass of notes is the example.
Timekeeping is the other function she singles out, because interrupting a person who talks
too long is a task human moderators dislike.
An agenda signal, including the news that the group is out of time and that the agenda was
badly planned, belongs with it.
She would also have the system restate, at the end, which decisions it has understood, as a
check rather than as a declaration.
A further use, which she distinguishes from moderation, is to invite an AI persona as a
customer or as a missing expert.
That persona is asked what it has heard and what concerns it.
E1 treats this as a different layer from human facilitation: hybrid teams in which agents
hold roles.
E2's list is procedural.
A warning that a named bias, such as the halo effect, may be operating would be useful, as
would a proposal to run a suitable decision procedure once several options are visible.
Enriching the option space, pressing for clarification, and prompting ``disagree and
commit'', ``safe enough to try'', or the cost of delay are further impulses he names.
He also says that the questions of preparation, including whether the decision owner is
present, may matter more than intervention in the live conversation, and that those
questions are easier to handle than real-time moderation.
E3 can imagine a decision mode that the group switches on: first the decision question,
including what happens if nothing is done, then one option at a time, then a visual summary
of advantages and disadvantages.
A second function is attentional.
The system reads the transcript and points at a place where a decision seems to be required,
or where the group is less agreed than it thinks.
Integration with the task system the team already uses is, for E3, part of a useful record.
E3 is explicit that the value of a skilled facilitator is not what this system would
reproduce.
The protocol and the guided decision are the parts that remain imaginable.

\begin{table}[htbp]
  \centering
  \caption{Limits the experts placed on an AI that intervenes in a meeting.}
  \label{tab:interview-ai-limits}
  \footnotesize
  \begin{tabularx}{\textwidth}{@{}>{\RaggedRight\arraybackslash}Xccc@{}}
    \toprule
    Limit & E1 & E2 & E3 \\
    \midrule
    Sentiment or attendance monitoring becomes surveillance & yes &     &     \\
    A record chills dissent, or leaves the dissenter marked & yes & yes &     \\
    Calling on a quiet person can shame, or can misread the silence & yes &     & yes \\
    People will not yield the floor to a machine, especially a senior person & yes & yes & yes \\
    The system must not decide, and must not impose a structure nobody asked for & yes & yes & yes \\
    A transcript is not enough context for a sound intervention & yes & yes & yes \\
    The group has to agree beforehand that interruption is allowed & yes &     & yes \\
    Replacing human talk removes the reason for meeting at all & yes &     &     \\
    \bottomrule
  \end{tabularx}
\end{table}

The limits constrain the functions.
E1 is willing to discuss sentiment analysis and immediately refuses it as a way of running
a meeting.
A system that can say a person has not spoken in the last ten meetings, or that can infer
mood from the voice, is, for her, a shift from contact to control, and it collides with data
protection.
The same record threatens the safe space in which critical comments are possible.
Those comments can later be used against the speaker.
Her concrete counterexample to a behavioral reading is a person who is only half present
because a child is ill.
The facilitator knows that because she was told.
The system does not, and it will still produce a score.
E2 describes the same chilling effect from the recording itself.
People become careful in the phase where reservations are most useful, and a reservation
that should have moved the discussion stays attached to the person who raised it.
He also doubts that a system without accepted authority will be obeyed.
E3's strongest objection is interruption.
People who want to speak will not let a machine take the floor.
The inhibition is higher with a human, including a beginner.
Silence, furthermore, is often not a trait and not a problem.
The person may be about to be persuaded.
An intervention at that moment can freeze them in the view they were leaving.
Being named by the system as the one who has not spoken can also shame them.
E3 would not exclude the possibility in principle.
The preconditions are a known mandate, video, tone, gesture, the history of this group, and
the meetings before and after this one.
A transcript does not supply them.
E2 makes a thinner version of the same point: a useful answer depends on the kind of
organization and on whether the work is remote or hybrid, and a general system that does
not know that context will stay vague.

How the system enters the meeting (A2) is where E1 and E3 are closest, and where they do
not say the same thing about initiative.
Both require an introduction.
E1 requires an agreement of the whole group that the system may interrupt, together with
an explicit decision about what is stored.
A check-in that contains personal information is, for her, a candidate for being kept out
of the record.
If the price of a clean record is that people stop speaking personally, she treats that
price as a loss of organizational culture.
She also distinguishes a moderator from a facilitator and would use only the first word
for the system.
A chat message is talked over, because people would rather talk to each other.
A video agent that tells a senior person they have spoken too long is, in her account,
likely to be rejected.
The human facilitator should be the one who calls the summary, or who knows the moment to
call it.
E3 goes further on initiative.
The group should trigger the decision mode, or the system should speak afterwards.
A light that says a decision still seems open, which the group can choose to follow, is
the most proactive form E3 will consider, and even that might damage the conversation.
An unsolicited structure that the participants did not ask for is not a form E3 can
imagine working.
E1, by contrast, does describe one unsolicited speech act as useful: the timekeeping
interruption, if the group has agreed to it.
The disagreement is about which acts may be proactive.
Summaries, documentation, and a decision procedure the group switches on are shared ground.
Taking the floor in order to correct a person is contested.

Two further boundaries belong in the result because they limit the claim an in-meeting
moderator can make.
E1 asks, of every proposed use, whether it produces a real gain or is a toy, and whether
it removes human interaction that the meeting existed to provide.
If the human exchange is no longer needed, the meeting itself is no longer needed.
E3 adds a practical bound: at some point a human facilitator is the cheaper instrument,
and operating breakouts and reading the arrangement of a shared board are part of what a
system would have to do if it claimed to run a longer decision process.
E2's image for the whole topic is an iceberg.
The meeting is the visible part.
Decision rights and the reasons people have to meet at all sit underneath.
Support, on that image, is justified where it can actually reach.
It is not justified by treating the meeting as the problem.

% ---------------------------------------------------------------------------------------------------------------------
\subsection{Relation to the challenges in Chapter~2}
\label{sec:interview-relation-to-challenges}

The interviews were not designed as a test of the literature, and they are not reported as
one.
Three contacts are close enough to state.

The failure of goal clarity and structure
(\autoref{chal:clear-goals-and-structure}) is present in all three accounts.
The experts add a location.
The unclear goal is often already true in the invitation and in the briefing, and the
structure that fails is often the decision procedure, not only the agenda of the hour.
Unequal participation and hierarchy
(\autoref{sec:key-challenges-in-meetings}) are present as well.
The remedy in these interviews is not equal speaking time.
It is a judgment about whose stake matters, combined with a way of speaking that does not
expose the person.
The groan zone and the diamond
(\autoref{sec:divergent-convergent}) are not only analytic labels here.
E1 describes the diamond as the shape of her workshops.
E2 names the groan zone and the duty not to close it under time pressure.
The neutrality of the facilitator
(\autoref{sec:facilitation}) matches E3's refusal to decide, and it matches E1's demand
that the group take ownership.
E1's consulting proposals show the boundary under strain.
She puts content in front of the group and still leaves the decision with them.

Two topics that Chapter~2 treats as central are only partly visible in the interviews.
The experts describe practices that seek missing perspectives and that delay evaluation
while positions are still being collected.
They do not describe the hidden-profile pattern in which the better option depends on
information that only one person holds.
E2 does propose that an AI might flag a named bias, including the halo effect.
That is a proposed function, not a report that the bias occurred in a meeting he facilitated.
Technical disturbance is outside the main argument of Chapter~2 and is still in the
material.
E3 names cross-organization tooling and ordinary audio failure as reasons a remote decision
meeting fails before facilitation starts.

What the interviews add, beyond the challenge list used to frame them, is the part of the
decision that an in-meeting moderator does not own.
That part includes who is allowed to decide, whether the participation is real, whether
the meeting is standing in for missing autonomy, whether the people affected are in the
room, and whether the result is owned and checked afterwards.
It also includes the remote loss of informal repair and of tone, and the ethical choice
about what a record is allowed to keep.
Those findings remain part of the result whether or not a later prototype implements them.
